\chapter{Guide d'entretien individuel de fin de projet}
\label{annexe:entretiens-fin-projet}

\section{Objectif de l'entretien}

À la fin du projet, un entretien individuel court peut être organisé avec chaque membre de l'équipe. Cet échange a pour objectif de recueillir son ressenti sur le déroulement du projet, sa place au sein du groupe, la mobilisation de ses compétences ainsi que les éventuelles difficultés rencontrées.

Il ne s'agit pas d'une évaluation individuelle, mais d'un temps de dialogue permettant de prendre du recul sur l'expérience collective. Les retours recueillis doivent contribuer à identifier les pratiques ayant bien fonctionné et les points qui pourraient être améliorés lors d'un futur projet.

\section{Organisation proposée}

\begin{itemize}
    \item \textbf{Durée indicative :} 20 à 25 minutes ;
    \item \textbf{Format :} entretien individuel semi-directif ;
    \item \textbf{Participants :} un membre de l'équipe et la personne chargée de la coordination du projet ;
    \item \textbf{Cadre :} échange bienveillant, confidentiel et sans jugement ;
    \item \textbf{Utilisation des réponses :} synthèse globale et anonymisée, sauf accord explicite de la personne interrogée.
\end{itemize}

Les questions proposées constituent une trame. Elles peuvent être adaptées selon les réponses de la personne et ne doivent pas nécessairement être posées dans un ordre strict.

\section{Introduction de l'entretien}

L'entretien peut être introduit de la manière suivante :

\begin{quote}
Cet échange a pour objectif de recueillir ton ressenti sur le projet et sur la manière dont nous avons travaillé ensemble. Il ne s'agit pas de t'évaluer, mais de comprendre ce qui a bien fonctionné, ce qui a été plus difficile et ce que nous pourrions améliorer pour un prochain projet. Tu peux répondre librement et illustrer tes réponses par des exemples.
\end{quote}

\section{Trame de questions}

\subsection*{1. Ressenti général}

\begin{quote}
Globalement, comment t'es-tu senti au cours du projet ?
\end{quote}

Relances possibles :

\begin{itemize}
    \item T'es-tu senti intégré et à ta place dans l'équipe ?
    \item Y a-t-il eu des moments où tu t'es senti en difficulté, en retrait ou en décalage ?
    \item Quels ont été les moments les plus satisfaisants ou les plus difficiles pour toi ?
\end{itemize}

\subsection*{2. Place et fonctionnement dans l'équipe}

\begin{quote}
Comment as-tu vécu ta place et ton rôle au sein de l'équipe ?
\end{quote}

Relances possibles :

\begin{itemize}
    \item Ton rôle et tes responsabilités étaient-ils suffisamment clairs ?
    \item As-tu eu suffisamment d'autonomie et de possibilités de participer aux décisions ?
    \item Comment as-tu vécu la communication et la collaboration entre les membres de l'équipe ?
\end{itemize}

\subsection*{3. Compétences et apprentissages}

\begin{quote}
Comment as-tu vécu la mobilisation de tes compétences pendant le projet ?
\end{quote}

Relances possibles :

\begin{itemize}
    \item As-tu eu l'occasion d'utiliser tes compétences à leur juste niveau ?
    \item Certaines de tes compétences auraient-elles pu être davantage mobilisées ?
    \item Qu'as-tu appris au cours du projet ?
    \item Sur quels sujets considères-tu avoir progressé ou gagné en autonomie ?
\end{itemize}

\subsection*{4. Difficultés et points de blocage}

\begin{quote}
Qu'est-ce qui t'a le plus aidé à avancer et, au contraire, qu'est-ce qui a pu te freiner ?
\end{quote}

Les éléments évoqués peuvent concerner l'organisation, les outils, les choix techniques, la communication, la charge de travail, les délais ou le fonctionnement collectif.

\subsection*{5. Gestion du projet et management}

\begin{quote}
Comment as-tu vécu la manière dont le projet et l'équipe ont été coordonnés ?
\end{quote}

Relances possibles :

\begin{itemize}
    \item Les consignes, les priorités et les décisions étaient-elles suffisamment claires ?
    \item T'es-tu senti suffisamment accompagné et écouté ?
    \item Le niveau d'autonomie proposé te convenait-il ?
    \item Aurais-tu préféré une organisation ou un style de management différent ?
    \item Y a-t-il une pratique de management qu'il faudrait conserver ou modifier ?
\end{itemize}

\subsection*{6. Améliorations possibles}

\begin{quote}
Selon toi, que pourrions-nous améliorer dans notre manière de travailler ?
\end{quote}

La réponse peut porter sur :

\begin{itemize}
    \item la gestion et la planification du projet ;
    \item la répartition des rôles et de la charge de travail ;
    \item la communication au sein de l'équipe ;
    \item les choix techniques et les méthodes de développement ;
    \item les outils utilisés ;
    \item les relations avec le client ou les parties prenantes.
\end{itemize}

\subsection*{7. Projection sur un futur projet}

\begin{quote}
Si le projet devait être recommencé, que conserverais-tu et que changerais-tu ?
\end{quote}

Relances possibles :

\begin{itemize}
    \item Que faudrait-il mettre en place plus tôt ?
    \item Quelle décision aurait pu être prise différemment ?
    \item Quel conseil donnerais-tu à une équipe qui reprendrait un projet similaire ?
\end{itemize}

\subsection*{8. Question ouverte}

\begin{quote}
Y a-t-il un sujet que nous n'avons pas abordé et que tu souhaiterais ajouter ?
\end{quote}

Cette dernière question permet à la personne de faire remonter un élément important qui ne serait pas apparu dans la trame initiale.

\section{Retour mutuel et conclusion}

Afin que l'entretien constitue réellement un temps d'échange, la personne chargée de la coordination peut également formuler un retour à destination du membre de l'équipe. Ce retour peut comprendre :

\begin{itemize}
    \item une contribution ou une qualité particulièrement appréciée pendant le projet ;
    \item une progression observée ;
    \item un axe de développement proposé de manière constructive.
\end{itemize}

L'entretien peut se terminer en résumant les principaux éléments exprimés afin de vérifier qu'ils ont été correctement compris :

\begin{quote}
Si je résume, les principaux points que tu retiens sont [...] et les améliorations que tu proposes sont [...]. Est-ce que cette synthèse correspond bien à ce que tu souhaitais exprimer ?
\end{quote}

\section{Exploitation des retours}

Les entretiens peuvent faire l'objet d'une synthèse transversale, sans attribuer nominativement les remarques. Cette synthèse peut notamment faire apparaître :

\begin{itemize}
    \item les points forts les plus fréquemment mentionnés ;
    \item les difficultés communes à plusieurs membres ;
    \item les compétences développées au cours du projet ;
    \item les propositions d'amélioration pour de futurs travaux collectifs ;
    \item les éventuels écarts de perception entre les membres de l'équipe.
\end{itemize}

Ces résultats peuvent alimenter la rétrospective du projet et contribuer à une analyse plus complète de son organisation, au-delà des seuls résultats techniques obtenus.
